\documentclass[UTF8,12pt,a4paper]{ctexart}
\usepackage[a4paper,margin=2.30cm,headheight=15pt]{geometry}
\usepackage{amsmath,amssymb,mathtools,bm}
\usepackage{booktabs,array,longtable}
\usepackage{xcolor,enumitem,fancyhdr,hyperref,listings}

\definecolor{sourcegray}{RGB}{75,84,97}
\definecolor{sourcebg}{RGB}{244,247,250}
\definecolor{linkblue}{RGB}{31,78,121}
\definecolor{keyblue}{RGB}{17,83,142}
\definecolor{keybg}{RGB}{239,247,255}
\hypersetup{colorlinks=true,linkcolor=linkblue,citecolor=linkblue,urlcolor=linkblue,
  pdftitle={第五题：投影法的有限体积离散与OpenFOAM求解器},pdfauthor={张云飞}}
\pagestyle{fancy}
\fancyhf{}
\lhead{第五题：不可压Navier--Stokes方程}
\rhead{投影法}
\cfoot{\thepage}
\setlength{\parindent}{2em}
\setlength{\parskip}{0.25em}
\setlist[itemize]{leftmargin=2.2em,itemsep=0.25em,topsep=0.35em}
\setlist[enumerate]{leftmargin=2.4em,itemsep=0.35em,topsep=0.35em}
\numberwithin{equation}{section}

\newcommand{\dd}{\,\mathrm d}
\newcommand{\dt}{\Delta t}
\newcommand{\Task}{\textnormal{[题目]}}
\newcommand{\Chorin}{\textnormal{[8]}}
\newcommand{\Guermond}{\textnormal{[9]}}
\newcommand{\Rone}{\textnormal{[R1]}}
\newcommand{\Rtwo}{\textnormal{[R2]}}
\newcommand{\Rthree}{\textnormal{[R3]}}
\newcommand{\OFfourteen}{\textnormal{[OF14]}}
\newcommand{\source}[1]{\par\smallskip\noindent\colorbox{sourcebg}{%
  \parbox{0.955\linewidth}{\footnotesize\color{sourcegray}\textbf{参考来源：}#1}}\par\smallskip}
\newcommand{\keytitle}[1]{\par\smallskip\noindent\colorbox{keybg}{%
  \parbox{0.965\linewidth}{\color{keyblue}\textbf{重要公式：}#1}}\par\smallskip}
\lstdefinestyle{foam}{
  basicstyle=\ttfamily\small,frame=single,framerule=0.4pt,
  rulecolor=\color{sourcegray},backgroundcolor=\color{sourcebg},
  numbers=left,numberstyle=\tiny\color{sourcegray},numbersep=8pt,
  columns=fullflexible,keepspaces=true,showstringspaces=false,
  breaklines=true,xleftmargin=1.4em,framexleftmargin=1.0em}

\title{\bfseries 第五题：不可压Navier--Stokes方程\\投影法、有限体积离散与OpenFOAM求解器详细推导}
\author{张云飞}
\date{2026年8月27日}

\begin{document}
\maketitle
\begin{abstract}
本文严格围绕题目第5.1节，依据Chorin的原始分步思想\Chorin 与Guermond、Minev、Shen对投影方法的系统分类\Guermond，推导常密度不可压Navier--Stokes方程的一阶非增量压力校正投影法。随后按照Moukalled、Mangani、Darwish的有限体积离散主线\Rone，将预测动量方程、压力Poisson方程、单元速度修正和面体积通量修正逐项离散到一般非正交混合非结构网格，并给出与OpenFOAM算子、线性方程和时间循环对应的求解器骨架。全文沿用前四题统一记号：内部面全局方向为$O_f\to N_f$，$\bm S_f=|S_f|\bm n_f$，局部装配方向由$\sigma_{cf}$给出；不使用$d(f)$表示相邻单元，也不使用$A_f$表示面面积。本文只推导投影法；PISO另见独立文档。
\end{abstract}
\tableofcontents
\newpage

\section*{题目范围、参考资料和统一记号}
\addcontentsline{toc}{section}{题目范围、参考资料和统一记号}
\subsection*{题目要求}
题目第5.1节给出
\begin{align}
  \nabla\cdot\bm u&=0,\label{eq:continuity}\\
  \frac{\partial\bm u}{\partial t}+\nabla\cdot(\bm u\otimes\bm u)
  &=-\frac{1}{\rho}\nabla p+\frac{1}{\rho}\nabla\cdot(\mu\nabla\bm u).
  \label{eq:momentum}
\end{align}
题目要求分别基于投影法和PISO方法写出求解过程，再使用有限体积法及OpenFOAM提供的离散格式开发数值求解器。本文件完成其中的投影法部分。
\source{\Task，第5.1节；投影法指定参考文献为\Chorin、\Guermond。}

\subsection*{常密度与运动学压力}
题目算例中的$\rho$和$\mu$取常数。定义
\begin{equation}
  \boxed{\nu:=\frac{\mu}{\rho},\qquad \pi:=\frac{p}{\rho}.}
  \label{eq:kinematic}
\end{equation}
于是
\begin{align}
  \nabla\cdot\bm u&=0,\label{eq:continuity_pi}\\
  \frac{\partial\bm u}{\partial t}+\nabla\cdot(\bm u\otimes\bm u)
  &=-\nabla\pi+\nabla\cdot(\nu\nabla\bm u).
  \label{eq:momentum_pi}
\end{align}
OpenFOAM不可压求解器中的字段\texttt{p}通常存储$\pi=p/\rho$，因此后文数学公式用$\pi$，代码仍使用字段名\texttt{p}。若$\rho=1$，二者数值相同，但量纲含义仍应区分。
\source{\Rone，第15章不可压流动控制方程及压力--速度耦合；\Rtwo，不可压OpenFOAM求解器中运动学压力的量纲约定。}

\subsection*{统一记号}
\begin{longtable}{>{\centering\arraybackslash}p{0.21\linewidth}p{0.68\linewidth}}
\toprule
记号 & 含义\\
\midrule
$c$ & 当前控制体\\
$\Omega_c,V_c$ & 控制体$c$及其体积（二维为面积乘单位厚度）\\
$\mathcal F_c$ & 控制体$c$的全部面集合\\
$f$ & 内部面\\
$O_f,N_f$ & 面$f$的owner与neighbour单元\\
$\bm d_f=\bm x_{N_f}-\bm x_{O_f}$，$d_f=\lVert\bm d_f\rVert$ & 从$O_f$到$N_f$的中心连线及距离\\
$\bm n_f$ & 从$O_f$指向$N_f$的全局单位法向量\\
$\bm S_f=|S_f|\bm n_f$ & 从$O_f$指向$N_f$的全局有向面积向量\\
$\sigma_{cf}$ & $c=O_f$时为$+1$，$c=N_f$时为$-1$\\
$\bm S_{cf}=\sigma_{cf}\bm S_f$ & 面$f$相对于控制体$c$的外向面积向量\\
$F_f=\bm u_f\cdot\bm S_f$ & 按$O_f\to N_f$定向的体积通量\\
$Q_c^m=\sum_f\sigma_{cf}F_f+Q_c^{m,\mathrm{bnd}}$ & 未除体积的离散通量不平衡\\
$R_c^m=Q_c^m/V_c$ & 体积归一化的离散连续性残差\\
$\bm H_f^{\mathrm{adv}},\bm H_f^{\mathrm{diff}}$ & 全局定向的对流、黏性动量通量\\
\bottomrule
\end{longtable}
\begin{equation}
  \boxed{\sigma_{cf}=\begin{cases}+1,&c=O_f,\\-1,&c=N_f,\end{cases}
  \qquad\bm S_{cf}=\sigma_{cf}\bm S_f,\qquad\bm S_f=|S_f|\bm n_f.}
  \label{eq:orientation}
\end{equation}
\source{\Rtwo，第1.3节：有限体积网格、面积向量、owner--neighbour寻址和内部面反号装配；与前四题统一记号一致。}

\subsection*{推导预备：为什么不可压方程需要投影}
\subsubsection*{压力约束与Helmholtz分解}
压力没有独立的演化方程，而是作为Lagrange乘子，使新速度落入无散空间
\[
  \mathcal V_0:=\{\bm v:\nabla\cdot\bm v=0\}.
\]
设预测速度分解为
\begin{equation}
  \bm u^*=\bm u^{n+1}+\dt\nabla\pi^{n+1}.
  \label{eq:helmholtz_scaled}
\end{equation}
于是
\begin{equation}
  \boxed{\bm u^{n+1}=\bm u^*-\dt\nabla\pi^{n+1}.}
  \label{eq:velocity_correction_cont}
\end{equation}
取散度并使用$\nabla\cdot\bm u^{n+1}=0$：
\begin{equation}
  \boxed{\nabla^2\pi^{n+1}=\frac{1}{\dt}\nabla\cdot\bm u^*.}
  \label{eq:pressure_poisson_cont}
\end{equation}
压力Poisson方程不是额外假设，而是由速度修正和无散约束直接推出。
\keytitle{投影法的核心闭环是式\eqref{eq:pressure_poisson_cont}与式\eqref{eq:velocity_correction_cont}。}
\source{\Chorin：中间速度、压力Poisson方程及速度修正；\Guermond，第2--3节：Helmholtz分解与非增量压力校正。}

\subsection*{本文采用的一阶非增量压力校正投影法}
\Guermond 将投影类方法分为压力校正、速度校正和一致分裂等类别。为与\Chorin 原始思想对应，本文选择一阶非增量压力校正，不混入增量式或旋转式。

\subsubsection*{连续层面的步骤一：预测速度}
\begin{equation}
  \boxed{\frac{\bm u^*-\bm u^n}{\dt}
  +\nabla\cdot(\bm u^n\otimes\bm u^*)
  -\nabla\cdot(\nu\nabla\bm u^*)=\bm0.}
  \label{eq:predictor_semiimplicit}
\end{equation}
对流输运速度取$\bm u^n$，被输运速度与黏性项取$\bm u^*$，形成线性半隐式预测方程。若全部取$n$层，则得到完全显式预测；投影步骤不变。半隐式形式便于用OpenFOAM的\texttt{fvm}算子构造稳定矩阵。
\source{\Chorin：无新压力的中间速度；\Guermond，第3节非增量压力校正；\Rone，第8、11、13章。半隐式线性化属于有限体积实现选择。}

\subsubsection*{连续层面的步骤二和三：压力与修正}
\begin{align}
  \boxed{\nabla^2\pi^{n+1}
  =\frac{1}{\dt}\nabla\cdot\bm u^*,}
  \label{eq:projection_poisson}\\
  \boxed{\bm u^{n+1}=\bm u^*-\dt\nabla\pi^{n+1}.}
  \label{eq:projection_corrector}
\end{align}
代回即有$\nabla\cdot\bm u^{n+1}=0$。

\subsubsection*{压力边界和参考值}
若$\bm u^{n+1}=\bm g^{n+1}$，由速度修正式的法向分量得到
\begin{equation}
  \boxed{\frac{\partial\pi^{n+1}}{\partial n}
  =\frac{1}{\dt}(\bm u^*-\bm g^{n+1})\cdot\bm n.}
  \label{eq:pressure_neumann}
\end{equation}
不可穿透壁面且预测速度满足正确法向速度时，右端为零。纯Neumann问题只确定到常数，必须规定
\begin{equation}
  \boxed{\pi_{c_{\mathrm{ref}}}=0.}
  \label{eq:pressure_reference}
\end{equation}
\source{\Guermond：投影边界条件与分裂误差；\Rone，第15.6.2.5节：压力相对性。}

\section{第一部分：控制方程的有限体积积分与半离散形式}
\subsection{从微分方程到控制体守恒式}
对任意控制体$\Omega_c$，连续性方程积分为
\begin{equation}
  \int_{\Omega_c}\nabla\cdot\bm u\,\mathrm dV=0.
\end{equation}
对左端使用Gauss定理，再用一个面中心积分点近似每个面的积分：
\begin{equation}
  \int_{\partial\Omega_c}\bm u\cdot\bm n\,\mathrm dS
  \approx\sum_{f\in\mathcal F_c}\bm u_f\cdot\bm S_{cf}=0.
\end{equation}
内部面采用全局方向$O_f\to N_f$后，控制体$c$的未除体积质量不平衡为
\begin{equation}
  \boxed{Q_c^m(\bm u):=
  \sum_{f\in\mathcal F_c^{\mathrm{int}}}\sigma_{cf}F_f
  +Q_c^{m,\mathrm{bnd}}=0,
  \qquad F_f:=\bm u_f\cdot\bm S_f.}
  \label{eq:semidiscrete_continuity}
\end{equation}
除以体积得到离散散度
\begin{equation}
  \boxed{R_c^m(\bm u):=\frac{Q_c^m(\bm u)}{V_c}=0.}
  \label{eq:semidiscrete_mass_residual}
\end{equation}

将运动量方程写成运动学形式
\begin{equation}
  \frac{\partial\bm u}{\partial t}
  +\nabla\cdot(\bm u\otimes\bm u)
  -\nabla\cdot(\nu\nabla\bm u)+\nabla\pi=\bm0,
\end{equation}
在$\Omega_c$上逐项积分。单元固定且采用单元中心值表示体积分时，
\begin{align}
  \int_{\Omega_c}\frac{\partial\bm u}{\partial t}\,\mathrm dV
  &\approx V_c\frac{\mathrm d\bm u_c}{\mathrm dt},\\
  \int_{\Omega_c}\nabla\cdot(\bm u\otimes\bm u)\,\mathrm dV
  &=\int_{\partial\Omega_c}(\bm u\otimes\bm u)\cdot\bm n\,\mathrm dS
  \approx\sum_f(\bm u_f\cdot\bm S_{cf})\bm u_f,\\
  \int_{\Omega_c}\nabla\cdot(\nu\nabla\bm u)\,\mathrm dV
  &\approx\sum_f\nu_f(\nabla\bm u)_f\cdot\bm S_{cf},\\
  \int_{\Omega_c}\nabla\pi\,\mathrm dV
  &\approx V_c(\nabla\pi)_c.
\end{align}
因此，在尚未指定任何面插值之前，有限体积半离散动量方程为
\begin{equation}
  \boxed{V_c\frac{\mathrm d\bm u_c}{\mathrm dt}
  +\sum_{f\in\mathcal F_c^{\mathrm{int}}}\sigma_{cf}\bm H_f^{\mathrm{adv}}
  -\sum_{f\in\mathcal F_c^{\mathrm{int}}}\sigma_{cf}\bm H_f^{\mathrm{diff}}
  +V_c(\nabla\pi)_c+\bm Q_c^{u,\mathrm{bnd}}=\bm0,}
  \label{eq:momentum_semidiscrete}
\end{equation}
其中
\begin{equation}
  \bm H_f^{\mathrm{adv}}:=F_f\bm u_f^{\mathrm{adv}},
  \qquad
  \bm H_f^{\mathrm{diff}}:=\nu_f(\nabla\bm u)_f\cdot\bm S_f.
\end{equation}
至此$F_f,\bm u_f^{\mathrm{adv}},(\nabla\bm u)_f$和$(\nabla\pi)_c$仍是未封闭的面量或梯度；它们必须在第二部分由空间离散给出，不能在控制体积分阶段暗中假定。

\subsection{半离散系统的约束结构}
把式\eqref{eq:momentum_semidiscrete}中除时间项外的全部项记为$\bm R_c^u(\bm u,\pi)$，则
\begin{equation}
  \boxed{V_c\frac{\mathrm d\bm u_c}{\mathrm dt}
  +\bm R_c^u(\bm u,\pi)=\bm0,
  \qquad Q_c^m(\bm u)=0.}
  \label{eq:semi_discrete_dae}
\end{equation}
第一式是演化方程，第二式是代数约束，因而不可压半离散系统是微分--代数系统。投影步骤的任务就是在每个时间步把预测通量送回$Q_c^m=0$的离散无散空间。
\keytitle{第一部分的最终结果是式\eqref{eq:semidiscrete_continuity}、\eqref{eq:momentum_semidiscrete}和\eqref{eq:semi_discrete_dae}；此时只完成守恒积分，尚未完成面插值和时间推进。}
\source{\Rone，第5章控制体积分与守恒离散，第8章扩散项，第11章对流项，第13章非稳态项，第15.1节不可压压力--速度约束。}

\section{第二部分：空间离散、面插值与代数装配}
\subsection{对流输运速度与守恒面通量}
对流通量需要区分两个量：$F_f$是跨面的体积通量，$\bm u_f^{\mathrm{adv}}$是被输运速度。前者决定方向并必须在相邻单元之间成对守恒；后者由对流格式插值。

一阶迎风可写为
\begin{equation}
  \bm u_f^{\mathrm{up}}=
  \begin{cases}
    \bm u_{O_f},&F_f\ge0,\\
    \bm u_{N_f},&F_f<0,
  \end{cases}
\end{equation}
或
\begin{equation}
  \boxed{\bm H_f^{\mathrm{adv}}
  =F_f^+\bm u_{O_f}+F_f^-\bm u_{N_f},}
  \label{eq:upwind_momentum_flux}
\end{equation}
其中$(F)^+=\max(F,0)$，$(F)^-=\min(F,0)$。若选\texttt{linearUpwind}，OpenFOAM用迎风值加梯度外推修正构造高阶面值，方向仍由$F_f$和$\sigma_{cf}$控制。
具体地，令$U_f$表示迎风单元，则
\begin{equation}
  \bm u_f^{\mathrm{LU}}
  =\bm u_{U_f}+(\nabla\bm u)_{U_f}\cdot(\bm x_f-\bm x_{U_f}),
\end{equation}
梯度由\texttt{gradSchemes}重构，必要时再施加限制器。无论选何种被输运量插值，内部面通量都只计算一次，并执行
\begin{equation}
  \bm Q_{O_f}^u\mathrel{+}=\bm H_f,
  \qquad
  \bm Q_{N_f}^u\mathrel{-}=\bm H_f.
\end{equation}
\source{\Rone，第11--12章；\Rtwo，OpenFOAM的\texttt{divSchemes}。}

\subsection{黏性面通量的正交与非正交分解}
\begin{equation}
  \bm S_f=\bm E_f+\bm T_f,\qquad \bm E_f\parallel\bm d_f,
  \label{eq:nonorth_decomp}
\end{equation}
则
\begin{align}
  \bm H_f^{\mathrm{diff}}
  &\approx D_f^u(\bm u_{N_f}-\bm u_{O_f})
  +\nu_f(\nabla\bm u)_f\cdot\bm T_f,
  \label{eq:diff_flux_nonorth}\\
  D_f^u&:=\nu_f\frac{|\bm E_f|}{d_f}.
\end{align}
第一项隐式进入矩阵，第二项显式修正，对应\texttt{Gauss linear corrected}。
\source{\Rone，第8章非正交扩散；\Rtwo，\texttt{laplacianSchemes}和\texttt{snGradSchemes}。}

\subsection{压力梯度的两种离散用途}
压力在求解器中以两种不同但必须一致的方式出现。动量修正需要单元中心梯度，Gauss线性格式给出
\begin{equation}
  \boxed{(\nabla\pi)_c\approx\frac{1}{V_c}
  \sum_{f\in\mathcal F_c}\pi_f\bm S_{cf},
  \qquad
  \pi_f\approx\lambda_f\pi_{O_f}+(1-\lambda_f)\pi_{N_f}.}
  \label{eq:cell_pressure_gradient}
\end{equation}
压力Poisson方程需要面法向梯度：
\begin{equation}
  \boxed{(\nabla\pi)_f\cdot\bm S_f
  \approx\frac{|\bm E_f|}{d_f}(\pi_{N_f}-\pi_{O_f})
  +(\nabla\pi)_f\cdot\bm T_f.}
  \label{eq:face_pressure_snGrad}
\end{equation}
前者对应\texttt{fvc::grad(p)}，后者对应\texttt{fvm::laplacian(coeff,p)}及该矩阵的\texttt{flux()}。速度修正和面通量修正不能各自另选一套压力梯度，否则压力方程求得的通量抵消关系会被破坏。
\source{\Rone，第9章梯度重构、第15.5.2节压力方程；\Rtwo，\texttt{gradSchemes}、\texttt{snGradSchemes}和\texttt{laplacianSchemes}。}

\subsection{由面贡献到单元矩阵}
把迎风对流与正交扩散的内部面贡献分别写到owner和neighbour方程中。以$F_f\ge0$为例，面$f$对owner方程提供$+F_f\bm u_{O_f}$，对neighbour方程提供$-F_f\bm u_{O_f}$；扩散项则产生成对的$D_f^u(\bm u_{O_f}-\bm u_{N_f})$。遍历所有面并加入边界贡献后，每个速度分量形成
\begin{equation}
  \boxed{V_c\frac{\mathrm d\bm u_c}{\mathrm dt}
  +a_c^{\mathrm{sp}}\bm u_c+
  \sum_{f\in\mathcal F_c^{\mathrm{int}}}
  a_{cf}^{\mathrm{sp}}\bm u_{\mathrm{other}(c,f)}
  =\bm b_c^{\mathrm{sp}}-V_c(\nabla\pi)_c.}
  \label{eq:predictor_algebra}
\end{equation}
“other”只用于文字说明，实际实现不定义$d(f)$：内部面直接以$O_f,N_f$成对装配。$a_c^{\mathrm{sp}}$和$a_{cf}^{\mathrm{sp}}$来自隐式对流、正交扩散与边界线性化；$\bm b_c^{\mathrm{sp}}$含显式非正交修正和显式源项。\texttt{fvm}算子负责左端矩阵，\texttt{fvc}算子只返回已知场。
\keytitle{第二部分的最终闭合关系是式\eqref{eq:upwind_momentum_flux}、\eqref{eq:diff_flux_nonorth}、\eqref{eq:cell_pressure_gradient}和\eqref{eq:face_pressure_snGrad}。尤其要区分守恒面通量$F_f$与被输运速度$\bm u_f^{\mathrm{adv}}$。}
\source{\Rone，第8、11、13章及第15.5.1节动量方程代数形式。}

\section{第三部分：时间离散与投影压力--速度耦合}
\subsection{从半离散动量方程得到预测方程}
在$t^n\to t^{n+1}$采用后向Euler，令输运通量取$n$层而被输运速度和黏性项取预测层。将式\eqref{eq:momentum_semidiscrete}中的新压力项暂时去掉，得到
\begin{equation}
  \boxed{\frac{V_c}{\dt}(\bm u_c^*-\bm u_c^n)
  +\sum_{f\in\mathcal F_c^{\mathrm{int}}}\sigma_{cf}\bm H_f^{\mathrm{adv},*}
  -\sum_{f\in\mathcal F_c^{\mathrm{int}}}\sigma_{cf}\bm H_f^{\mathrm{diff},*}
  +\bm Q_c^{u,\mathrm{bnd},*}=\bm0.}
  \label{eq:predictor_fvm_global}
\end{equation}
其中$\bm H_f^{\mathrm{adv},*}=F_f^n\bm u_f^*$，$F_f^n=\bm u_f^n\cdot\bm S_f$。时间项移项后，式\eqref{eq:predictor_fvm_global}形成
\begin{equation}
  \boxed{a_c^u\bm u_c^*+
  \sum_{f\in\mathcal F_c^{\mathrm{int}}}a_{cf}^u
  \bm u_{\mathrm{other}(c,f)}^*=\bm b_c^u,}
  \label{eq:projection_predictor_matrix}
\end{equation}
其中$a_c^u$包含$V_c/\dt$，$\bm b_c^u$包含$(V_c/\dt)\bm u_c^n$、显式非正交项和边界项。求解该线性系统得到$\bm u^*$，再用一致面插值构造$F_f^*=\bm u_f^*\cdot\bm S_f$。由于预测式没有施加新时刻连续性，通常$Q_c^m(\bm u^*)\ne0$。
\source{\Chorin；\Guermond，第3节非增量压力校正；\Rone，第13章时间离散和第15章压力--速度耦合。}

\subsection{从离散连续性逐步推出压力Poisson方程}
定义预测面通量$F_f^*:=\bm u_f^*\cdot\bm S_f$。为使离散连续性严格成立，直接修正面通量：
\begin{equation}
  \boxed{F_f^{n+1}
  =F_f^*-\dt(\nabla\pi^{n+1})_f\cdot\bm S_f.}
  \label{eq:flux_correction}
\end{equation}
对单元$c$要求
\begin{equation}
  Q_c^{m,n+1}:=
  \sum_{f\in\mathcal F_c^{\mathrm{int}}}\sigma_{cf}F_f^{n+1}
  +Q_c^{m,\mathrm{bnd},n+1}=0.
  \label{eq:discrete_continuity}
\end{equation}
代入得
\begin{equation}
  \boxed{\sum_{f\in\mathcal F_c}
  (\nabla\pi^{n+1})_f\cdot\bm S_{cf}
  =\frac{1}{\dt}\sum_{f\in\mathcal F_c}F_{cf}^*.}
  \label{eq:pressure_fvm}
\end{equation}
内部面$F_{cf}^*=\sigma_{cf}F_f^*$。这正是连续Poisson方程的守恒有限体积形式。
\keytitle{有限体积投影必须修正面通量，再逐控制体检查通量和。只修正单元速度后再插值到面，不能保证离散连续性严格成立。}
\source{\Chorin、\Guermond：Poisson方程由无散约束产生；\Rone，第15.5.2节：预测面通量代入连续性构造压力方程。}

\subsection{压力方程的正交与非正交离散}
正交网格上
\begin{equation}
  (\nabla\pi)_f\cdot\bm S_f
  \approx\frac{|S_f|}{d_f}(\pi_{N_f}-\pi_{O_f}).
\end{equation}
非正交网格上
\begin{equation}
  (\nabla\pi)_f\cdot\bm S_f
  \approx\frac{|\bm E_f|}{d_f}(\pi_{N_f}-\pi_{O_f})
  +(\nabla\pi)_f\cdot\bm T_f.
  \label{eq:p_nonorth}
\end{equation}
第一项隐式入矩阵，第二项在非正交校正循环中显式更新。只有最后一次校正将压力方程通量用于更新$F_f^{n+1}$。
将式\eqref{eq:p_nonorth}乘以$\dt$，定义内部面压力传递系数
\begin{equation}
  D_f^p:=\dt\frac{|\bm E_f|}{d_f}.
\end{equation}
忽略边界和显式非正交项时，面$f$对owner压力方程产生
\begin{equation}
  D_f^p(\pi_{N_f}^{n+1}-\pi_{O_f}^{n+1})=F_f^*,
\end{equation}
对neighbour方程符号相反。因此逐面装配后得到
\begin{equation}
  \boxed{a_c^p\pi_c^{n+1}
  +\sum_{f\in\mathcal F_c^{\mathrm{int}}}a_{cf}^p
  \pi_{\mathrm{other}(c,f)}^{n+1}=b_c^p,}
  \label{eq:projection_pressure_matrix}
\end{equation}
其中内部面的非对角系数由$-D_f^p$给出，对角系数为相邻传递系数之和，$b_c^p$由预测通量不平衡、边界通量以及显式非正交修正组成。这个装配过程正是\texttt{fvm::laplacian(dt,p)}所建立的矩阵。
\source{\Rone，第8章和第15.5.2--15.5.3节；\Rtwo，\texttt{corrected}面法向梯度。}

\subsection{单元速度与守恒面通量的同步修正}
\begin{align}
  \boxed{\bm u_c^{n+1}
  =\bm u_c^*-\dt(\nabla\pi^{n+1})_c,}
  \label{eq:cell_velocity_correction}\\
  \boxed{F_f^{n+1}=F_f^*-H_f^p,\qquad
  H_f^p:=\dt(\nabla\pi^{n+1})_f\cdot\bm S_f.}
  \label{eq:final_phi}
\end{align}
单元梯度可用
\begin{equation}
  (\nabla\pi)_c\approx \frac{1}{V_c}
  \sum_{f\in\mathcal F_c}\pi_f\bm S_{cf}.
\end{equation}
$H_f^p$必须采用与压力方程相同的正交/非正交离散。

\subsection{离散连续性的代数证明}
压力方程等价于
\[
  \sum_f\dt(\nabla\pi^{n+1})_f\cdot\bm S_{cf}
  =\sum_fF_{cf}^*.
\]
故
\begin{align}
  Q_c^{m,n+1}
  &=\sum_f\left[F_{cf}^*
  -\dt(\nabla\pi^{n+1})_f\cdot\bm S_{cf}\right]=0,\\
  \boxed{R_c^{m,n+1}:=\frac{Q_c^{m,n+1}}{V_c}=0.}
  \label{eq:discrete_zero_div}
\end{align}
实际“零”由压力线性求解器容差决定。

\subsection{投影边界条件、压力零空间与时间精度}
方腔顶盖取$\bm u=(1,0)$，其余壁面$\bm u=\bm0$；三角腔顶边取$\bm u=(1,0)$，两条斜壁$\bm u=\bm0$。速度均为Dirichlet条件，不可穿透壁面满足$F_b=0$。

若预测步骤保持正确壁面法向速度，则壁面运动学压力取$\partial\pi/\partial n=0$，OpenFOAM中为\texttt{zeroGradient}。所有速度法向分量均指定时，压力矩阵有常数零空间，必须设置\texttt{pRefCell}和\texttt{pRefValue}。
本文采用的非增量投影在一般边界条件下具有一阶时间精度并存在分裂误差；若改用增量式或旋转式，压力变量、右端和边界处理都必须随之改变，不能只替换一个更新公式。
\source{\Rone，第15.6节；\Guermond：投影法边界条件。}

\section{第四部分：OpenFOAM求解器构造、配置与验证}
\subsection{开发路线：新增求解器模块而不是虚构控制器}
OpenFOAM~14的通用可执行程序是\texttt{foamRun}。它从\texttt{controlDict}读取
\texttt{solver}并动态加载求解器模块。单个时间步内，驱动程序先调用预处理和动量预测，
再调用压力校正与步末处理。官方\texttt{incompressibleFluid}模块实现PIMPLE/PISO，
并没有名为\texttt{projection}的控制器。因此，投影法的正确开发路线是：
以\texttt{incompressibleFluid}为基类建立运行时可选模块\texttt{projectionFluid}，
复用其网格、$\bm u$、$\pi$、$F_f$、输运模型、\texttt{fvModels}与
\texttt{fvConstraints}，仅重写无压力动量预测和投影压力校正。

\begin{equation}
  \boxed{\begin{aligned}
  \texttt{foamRun}&\longrightarrow
  \texttt{projectionFluid::momentumPredictor()}\\
  &\longrightarrow
  \texttt{projectionFluid::pressureCorrector()}.
  \end{aligned}}
  \label{eq:of14_lifecycle}
\end{equation}
\source{\OFfourteen：\texttt{foamRun.C}第109--193行的模块实例化与求解生命周期；官方
\texttt{incompressibleFluid}类的字段和虚函数接口。}

\subsection{状态、临时量和线性系统的生命周期}
\begin{longtable}{p{0.20\linewidth}p{0.28\linewidth}p{0.41\linewidth}}
\toprule
对象 & OpenFOAM类型 & 生命周期和职责\\
\midrule
$\bm u^n,\bm u^{n+1}$ & \texttt{volVectorField U\_} & 持久状态；读入、写出，并在时间层间保存\\
$\pi^{n+1}$ & \texttt{volScalarField p\_} & 持久运动学压力；纯Neumann时由\texttt{pressureReference}固定常数\\
$F_f^{n+1}$ & \texttt{surfaceScalarField phi\_} & 持久守恒面通量；压力校正后不得由\texttt{fvc::flux(U)}覆盖\\
$\bm u^*$ & \texttt{volVectorField UStar\_} & 模块成员；预测阶段求解，投影完成后同步为接受速度\\
$F_f^*$ & \texttt{surfaceScalarField phiStar\_} & 模块成员；由预测速度构造并满足边界通量约束\\
预测矩阵 & \texttt{tmp<fvVectorMatrix> tUStarEqn\_} & 在预测阶段装配、约束、求解，投影结束后可清除\\
压力矩阵 & \texttt{fvScalarMatrix pEqn} & 每个非正交校正中重新装配；最后一次才更新\texttt{phi\_}\\
\bottomrule
\end{longtable}
这一区分是求解器实现的关键：场保存算法状态，\texttt{fvMatrix}只保存当前线性化系统；
压力方程真正约束的是面通量，而不是一个事后插值得到的速度散度。

\subsection{数学公式与实际算子逐项对应}
\begin{longtable}{p{0.35\linewidth}p{0.53\linewidth}}
\toprule
数学或离散对象 & OpenFOAM~14表达\\
\midrule
$\partial_t\bm u^*$ & \texttt{fvm::ddt(UStar\_)}\\
$\nabla\cdot(\bm u^n\otimes\bm u^*)$ & \texttt{fvm::div(phi\_, UStar\_)}\\
黏性/动量输运 & \texttt{momentumTransport->divDevSigma(UStar\_)}\\
体源与矩阵约束 & \texttt{fvModels().source(UStar\_)}, \texttt{fvConstraints().constrain(...)}\\
$F_f^*$ & \texttt{phiStar\_ = fvc::flux(UStar\_)}并施加边界通量约束\\
$\nabla\cdot(\dt\nabla\pi)$ & \texttt{fvm::laplacian(dtCoeff,p\_)}\\
$\nabla\cdot F^*$ & \texttt{fvc::div(phiStar\_)}\\
$H_f^p=\dt\nabla\pi\cdot\bm S_f$ & \texttt{pEqn.flux()}，与压力矩阵使用同一面离散\\
$F_f^{n+1}=F_f^*-H_f^p$ & \texttt{phi\_ = phiStar\_ - pEqn.flux()}\\
$\bm u^{n+1}=\bm u^*-\dt\nabla\pi$ & \texttt{U\_ = UStar\_ - dt*fvc::grad(p\_)}\\
\bottomrule
\end{longtable}
固定黏度层流时，\texttt{divDevSigma}退化到相应黏性应力离散；保留模型接口比在代码中
硬写\texttt{-fvm::laplacian(nu,UStar)}更符合OpenFOAM~14的模块开发方式。
\source{\Rone，第8、11、13、15章；\OFfourteen：官方
\texttt{incompressibleFluid/momentumPredictor.C}第36--54行的矩阵装配、松弛与约束方式。}

\subsection{两个核心成员函数的开发骨架}
下面不是独立的旧式\texttt{main()}，而是供\texttt{foamRun}调用的模块成员函数。为聚焦本题，
代码假定静止网格；若加入动网格，还必须照官方压力校正加入\texttt{Uf}、绝对/相对通量转换和
\texttt{fvc::correctUf}。

\begin{lstlisting}[style=foam,language=C++]
void projectionFluid::momentumPredictor()
{
    // At the start of a new time step U_ and phi_ are the accepted fields.
    UStar_ = U_;

    tUStarEqn_ =
    (
        fvm::ddt(UStar_)
      + fvm::div(phi_, UStar_)
      + MRF.DDt(UStar_)
      + momentumTransport->divDevSigma(UStar_)
     == fvModels().source(UStar_)
    );

    fvVectorMatrix& UStarEqn = tUStarEqn_.ref();
    UStarEqn.relax();
    fvConstraints().constrain(UStarEqn);
    UStarEqn.solve();                 // no -fvc::grad(p_) here
    UStar_.correctBoundaryConditions();
    fvConstraints().constrain(UStar_);

    phiStar_ = fvc::flux(UStar_);
    MRF.makeRelative(phiStar_);
    if (p_.needReference())
    {
        adjustPhi(phiStar_, UStar_, p_);
    }
}
\end{lstlisting}

\begin{lstlisting}[style=foam,language=C++]
void projectionFluid::pressureCorrector()
{
    volScalarField dtCoeff
    (
        IOobject
        (
            "dtCoeff", runTime.name(), mesh,
            IOobject::NO_READ, IOobject::NO_WRITE, false
        ),
        mesh,
        runTime.deltaT()
    );

    // Make pressure BCs consistent with UStar_, phiStar_ and dtCoeff.
    constrainPressure(p_, UStar_, phiStar_, dtCoeff, MRF);

    while (pimple.correctNonOrthogonal())
    {
        fvScalarMatrix pEqn
        (
            fvm::laplacian(dtCoeff, p_)
         == fvc::div(phiStar_)
        );
        pEqn.setReference
        (
            pressureReference.refCell(),
            pressureReference.refValue()
        );
        pEqn.solve();

        if (pimple.finalNonOrthogonalIter())
        {
            phi_ = phiStar_ - pEqn.flux();
        }
    }

    U_ = UStar_ - dtCoeff*fvc::grad(p_);
    U_.correctBoundaryConditions();
    fvConstraints().constrain(U_);
    UStar_ = U_;                     // accepted state for next time level
    continuityErrors();
    tUStarEqn_.clear();
}
\end{lstlisting}
这里的\texttt{pimple.correctNonOrthogonal()}只是复用OpenFOAM现有的非正交循环控制；
它不把投影法变成PISO。\texttt{fvSolution}中应令\texttt{nOuterCorrectors=1}、
\texttt{nCorrectors=1}，使每个时间步只执行一次预测和一次投影。

\subsection{必须保留的边界和守恒处理}
\begin{enumerate}
  \item \texttt{adjustPhi}只在压力需要参考值的封闭域等情形修正预测边界通量平衡。
  \item \texttt{constrainPressure}使压力边界的法向梯度与预测速度、预测通量及$\dt$系数相容。
  \item \texttt{pEqn.flux()}必须在最后一个非正交循环中直接写入\texttt{phi\_}；其离散与压力矩阵完全一致。
  \item 速度修正后依次执行速度边界修正和\texttt{fvConstraints}场约束。
  \item 固定网格时以上流程闭合；动网格时还要照官方\texttt{correctPressure.C}处理MRF、\texttt{Uf}和相对通量。
\end{enumerate}
\keytitle{开发层面的不可替换公式链为：先得到\texttt{phiStar\_}，再求
$\texttt{laplacian(dtCoeff,p\_)}=\texttt{div(phiStar\_)}$，最后同时更新
\texttt{phi\_}与\texttt{U\_}；二者必须来自同一个压力解。}
\source{\OFfourteen：官方\texttt{correctPressure.C}第47--124行中的
\texttt{adjustPhi}、\texttt{constrainPressure}、非正交循环、\texttt{pEqn.flux()}、
速度边界修正和\texttt{fvConstraints}调用次序；本节仅把动量响应系数由\texttt{rAtU}替换为
投影法严格推出的$\dt$。}

\subsection{模块文件结构与构建入口}
\begin{lstlisting}[style=foam]
projectionFluid/
|-- projectionFluid.H
|-- projectionFluid.C
|-- momentumPredictor.C
|-- pressureCorrector.C
`-- Make/
    |-- files
    `-- options
\end{lstlisting}
\texttt{projectionFluid.H}声明运行时类型、三个算法成员
\texttt{UStar\_}、\texttt{phiStar\_}、\texttt{tUStarEqn\_}以及两个重写函数；
实现文件用\texttt{addToRunTimeSelectionTable}注册模块。\texttt{Make/files}生成用户动态库
\path{$(FOAM_USER_LIBBIN)/libprojectionFluid.so}。\texttt{Make/options}应从本机
\path{$FOAM_MODULES/incompressibleFluid/Make/options}复制相同依赖，再加入基类模块链接。
构建与运行顺序为
\begin{lstlisting}[style=foam]
wmake libso projectionFluid
foamRun -solver projectionFluid
\end{lstlisting}
不在推导文档中臆造随版本变化的完整链接选项；实际开发时以安装的OpenFOAM~14源码为唯一接口基准。

\subsection{与PISO压力方程的区别}
PISO压力方程系数是动量矩阵对角倒数$r_A=1/\texttt{UEqn.A()}$，本投影法的修正系数是$\dt$：
\[
  \boxed{\nabla\cdot(\dt\nabla\pi)=\nabla\cdot\bm u^*.}
\]
因此不能为了复用官方\texttt{correctPressure.C}而把\texttt{dtCoeff}机械替换为\texttt{rAU}；
否则实现的已是动量矩阵压力校正而非本文的Chorin型投影。

\subsection{离散格式、控制字典与线性求解配置}
\begin{lstlisting}[style=foam]
ddtSchemes
{
    default         Euler;
}
gradSchemes
{
    default         Gauss linear;
    grad(UStar)     cellLimited Gauss linear 1;
}
divSchemes
{
    default         none;
    div(phi,UStar)  bounded Gauss linearUpwind grad(UStar);
}
laplacianSchemes
{
    default         Gauss linear corrected;
}
interpolationSchemes
{
    default         linear;
}
snGradSchemes
{
    default         corrected;
}
fluxRequired
{
    default         no;
    p;
}
\end{lstlisting}

\begin{lstlisting}[style=foam]
solvers
{
    p
    {
        solver          GAMG;
        smoother        GaussSeidel;
        tolerance       1e-8;
        relTol          0.01;
    }
    UStar
    {
        solver          smoothSolver;
        smoother        symGaussSeidel;
        tolerance       1e-8;
        relTol          0.01;
    }
}
PIMPLE
{
    nOuterCorrectors         1;
    nCorrectors              1;
    nNonOrthogonalCorrectors 1;
    momentumPredictor        yes;
}
\end{lstlisting}
\begin{lstlisting}[style=foam]
application     foamRun;
solver          projectionFluid;
\end{lstlisting}
压力参考由OpenFOAM~14的\texttt{pressureReference}读取相应求解控制；具体字典位置以本机模块基类为准。
容差和修正次数是算例配置，不是\Chorin 或\Guermond 规定的普适常数，应通过连续性残差和网格质量验证。
\source{\Rtwo、\Rthree：空间离散、线性求解和运行字典；\OFfourteen：\texttt{foamRun}和
\texttt{incompressibleFluid}模块控制结构。}

\subsection{完整算法与实现检查}
给定$\bm u^0$、$\pi^0$和初始面通量，对$n=0,1,\ldots$执行：
\begin{enumerate}
  \item \texttt{foamRun}推进时间并进入投影模块的动量预测成员函数。
  \item 将接受场\texttt{U\_}复制到\texttt{UStar\_}，用持久守恒通量\texttt{phi\_}装配
  预测矩阵；依次进行矩阵松弛、矩阵约束、线性求解、速度边界修正和场约束。
  \item 由$\bm u^*$构造\texttt{phiStar\_}，执行MRF/边界净通量处理。
  \item 驱动程序进入压力校正成员函数，创建值为$\dt$的\texttt{dtCoeff}，
  然后施加压力边界一致性约束。
  \item 在\texttt{pimple.correctNonOrthogonal()}循环内装配压力方程\eqref{eq:pressure_fvm}，
  设置参考压力并求解；只在最后一次用\texttt{pEqn.flux()}执行式\eqref{eq:final_phi}。
  \item 使用式\eqref{eq:cell_velocity_correction}更新\texttt{U\_}，重施速度边界和
  \texttt{fvConstraints}，但保留压力方程产生的\texttt{phi\_}。
  \item 计算$Q_c^{m,n+1}$和$R_c^{m,n+1}$，确认式\eqref{eq:discrete_zero_div}在容差内成立，
  然后清除预测矩阵并写出接受场。
  \item 重复时间推进，直至速度场、主涡位置、残差和连续性误差均稳定。
\end{enumerate}

\keytitle{投影法有限体积求解器的四个核心式为：预测动量方程\eqref{eq:predictor_fvm_global}、压力方程\eqref{eq:pressure_fvm}、面通量修正式\eqref{eq:final_phi}和单元速度修正式\eqref{eq:cell_velocity_correction}。}

\begin{enumerate}
  \item 每个内部面只计算一次，owner加、neighbour减。
  \item 数学物理压力$p$与代码运动学压力$\pi=p/\rho$不得混淆。
  \item 压力方程修正面通量，不能再用简单插值覆盖守恒通量。
  \item 纯Neumann压力边界必须设置参考单元和值。
  \item 预测速度法向边界与压力Neumann条件必须满足式\eqref{eq:pressure_neumann}。
  \item 逐单元检查$R_c^m$，并检查全局边界净通量。
  \item 用中心线速度、主涡中心、流线和参考数据验证，不能只看残差。
\end{enumerate}

\subsection{参考来源与使用范围}
\begin{longtable}{p{0.13\linewidth}p{0.76\linewidth}}
\toprule
编号 & 本文实际使用内容\\
\midrule
\Task & 第5.1节控制方程和任务要求。\\
\Chorin & 无压力中间速度、压力Poisson方程、速度投影修正。\\
\Guermond & 投影方法分类；非增量压力校正；Helmholtz投影、边界和分裂误差。\\
\Rone & 有限体积积分、扩散、梯度、对流、时间、压力方程与面通量。\\
\Rtwo & OpenFOAM网格数据结构、owner--neighbour、面积向量、字段和离散格式。\\
\Rthree & OpenFOAM字典、运行和求解器设置。\\
\OFfourteen & OpenFOAM~14的\texttt{foamRun}生命周期、\texttt{incompressibleFluid}字段接口、
动量矩阵装配和压力校正调用次序；投影模块在此架构上作最小算法替换。\\
\bottomrule
\end{longtable}

\begin{thebibliography}{99}
\bibitem{Chorin1967}
A.~J. Chorin,
“A numerical method for solving incompressible viscous flow problems,”
\emph{Journal of Computational Physics}, vol.~2, pp.~12--26, 1967.
\bibitem{Guermond2006}
J.-L. Guermond, P. Minev, and J. Shen,
“An overview of projection methods for incompressible flows,”
\emph{Computer Methods in Applied Mechanics and Engineering},
vol.~195, pp.~6011--6045, 2006.
\bibitem{Moukalled2016}
F. Moukalled, L. Mangani, and M. Darwish,
\emph{The Finite Volume Method in Computational Fluid Dynamics: An Advanced Introduction with OpenFOAM and Matlab},
Springer, 2016. 重点使用第5、8、9、11、13、15章。
\bibitem{Primer2021}
T. Holzmann and S. Rooney, \emph{The OpenFOAM Technology Primer}, 2021.
\bibitem{Holzinger2020}
G. Holzinger, \emph{OpenFOAM: A Little User-Manual}, 2020.
\bibitem{OpenFOAM14}
OpenFOAM Foundation, \emph{OpenFOAM~14 Source Code Guide}: 
\texttt{foamRun.C}, \texttt{applications/modules/incompressibleFluid/momentumPredictor.C},
\texttt{correctPressure.C} and class interface, 2026.\\
\url{https://cpp.openfoam.org/v14/foamRun_8C_source.html}\\
\url{https://cpp.openfoam.org/v14/incompressibleFluid_2momentumPredictor_8C_source.html}\\
\url{https://cpp.openfoam.org/v14/incompressibleFluid_2correctPressure_8C_source.html}
\end{thebibliography}
\end{document}
